% Options for packages loaded elsewhere
\PassOptionsToPackage{unicode}{hyperref}
\PassOptionsToPackage{hyphens}{url}
\documentclass[
  ignorenonframetext,
]{beamer}
\newif\ifbibliography
\usepackage{pgfpages}
\setbeamertemplate{caption}[numbered]
\setbeamertemplate{caption label separator}{: }
\setbeamercolor{caption name}{fg=normal text.fg}
\beamertemplatenavigationsymbolsempty
% remove section numbering
\setbeamertemplate{part page}{
  \centering
  \begin{beamercolorbox}[sep=16pt,center]{part title}
    \usebeamerfont{part title}\insertpart\par
  \end{beamercolorbox}
}
\setbeamertemplate{section page}{
  \centering
  \begin{beamercolorbox}[sep=12pt,center]{section title}
    \usebeamerfont{section title}\insertsection\par
  \end{beamercolorbox}
}
\setbeamertemplate{subsection page}{
  \centering
  \begin{beamercolorbox}[sep=8pt,center]{subsection title}
    \usebeamerfont{subsection title}\insertsubsection\par
  \end{beamercolorbox}
}
% Prevent slide breaks in the middle of a paragraph
\widowpenalties 1 10000
\raggedbottom
\AtBeginPart{
  \frame{\partpage}
}
\AtBeginSection{
  \ifbibliography
  \else
    \frame{\sectionpage}
  \fi
}
\AtBeginSubsection{
  \frame{\subsectionpage}
}
\usepackage{iftex}
\ifPDFTeX
  \usepackage[T1]{fontenc}
  \usepackage[utf8]{inputenc}
  \usepackage{textcomp} % provide euro and other symbols
\else % if luatex or xetex
  \usepackage{unicode-math} % this also loads fontspec
  \defaultfontfeatures{Scale=MatchLowercase}
  \defaultfontfeatures[\rmfamily]{Ligatures=TeX,Scale=1}
\fi
\usepackage{lmodern}
\usetheme[]{Montpellier}
\ifPDFTeX\else
  % xetex/luatex font selection
\fi
% Use upquote if available, for straight quotes in verbatim environments
\IfFileExists{upquote.sty}{\usepackage{upquote}}{}
\IfFileExists{microtype.sty}{% use microtype if available
  \usepackage[]{microtype}
  \UseMicrotypeSet[protrusion]{basicmath} % disable protrusion for tt fonts
}{}
\makeatletter
\@ifundefined{KOMAClassName}{% if non-KOMA class
  \IfFileExists{parskip.sty}{%
    \usepackage{parskip}
  }{% else
    \setlength{\parindent}{0pt}
    \setlength{\parskip}{6pt plus 2pt minus 1pt}}
}{% if KOMA class
  \KOMAoptions{parskip=half}}
\makeatother
\usepackage{color}
\usepackage{fancyvrb}
\newcommand{\VerbBar}{|}
\newcommand{\VERB}{\Verb[commandchars=\\\{\}]}
\DefineVerbatimEnvironment{Highlighting}{Verbatim}{commandchars=\\\{\}}
% Add ',fontsize=\small' for more characters per line
\usepackage{framed}
\definecolor{shadecolor}{RGB}{248,248,248}
\newenvironment{Shaded}{\begin{snugshade}}{\end{snugshade}}
\newcommand{\AlertTok}[1]{\textcolor[rgb]{0.94,0.16,0.16}{#1}}
\newcommand{\AnnotationTok}[1]{\textcolor[rgb]{0.56,0.35,0.01}{\textbf{\textit{#1}}}}
\newcommand{\AttributeTok}[1]{\textcolor[rgb]{0.13,0.29,0.53}{#1}}
\newcommand{\BaseNTok}[1]{\textcolor[rgb]{0.00,0.00,0.81}{#1}}
\newcommand{\BuiltInTok}[1]{#1}
\newcommand{\CharTok}[1]{\textcolor[rgb]{0.31,0.60,0.02}{#1}}
\newcommand{\CommentTok}[1]{\textcolor[rgb]{0.56,0.35,0.01}{\textit{#1}}}
\newcommand{\CommentVarTok}[1]{\textcolor[rgb]{0.56,0.35,0.01}{\textbf{\textit{#1}}}}
\newcommand{\ConstantTok}[1]{\textcolor[rgb]{0.56,0.35,0.01}{#1}}
\newcommand{\ControlFlowTok}[1]{\textcolor[rgb]{0.13,0.29,0.53}{\textbf{#1}}}
\newcommand{\DataTypeTok}[1]{\textcolor[rgb]{0.13,0.29,0.53}{#1}}
\newcommand{\DecValTok}[1]{\textcolor[rgb]{0.00,0.00,0.81}{#1}}
\newcommand{\DocumentationTok}[1]{\textcolor[rgb]{0.56,0.35,0.01}{\textbf{\textit{#1}}}}
\newcommand{\ErrorTok}[1]{\textcolor[rgb]{0.64,0.00,0.00}{\textbf{#1}}}
\newcommand{\ExtensionTok}[1]{#1}
\newcommand{\FloatTok}[1]{\textcolor[rgb]{0.00,0.00,0.81}{#1}}
\newcommand{\FunctionTok}[1]{\textcolor[rgb]{0.13,0.29,0.53}{\textbf{#1}}}
\newcommand{\ImportTok}[1]{#1}
\newcommand{\InformationTok}[1]{\textcolor[rgb]{0.56,0.35,0.01}{\textbf{\textit{#1}}}}
\newcommand{\KeywordTok}[1]{\textcolor[rgb]{0.13,0.29,0.53}{\textbf{#1}}}
\newcommand{\NormalTok}[1]{#1}
\newcommand{\OperatorTok}[1]{\textcolor[rgb]{0.81,0.36,0.00}{\textbf{#1}}}
\newcommand{\OtherTok}[1]{\textcolor[rgb]{0.56,0.35,0.01}{#1}}
\newcommand{\PreprocessorTok}[1]{\textcolor[rgb]{0.56,0.35,0.01}{\textit{#1}}}
\newcommand{\RegionMarkerTok}[1]{#1}
\newcommand{\SpecialCharTok}[1]{\textcolor[rgb]{0.81,0.36,0.00}{\textbf{#1}}}
\newcommand{\SpecialStringTok}[1]{\textcolor[rgb]{0.31,0.60,0.02}{#1}}
\newcommand{\StringTok}[1]{\textcolor[rgb]{0.31,0.60,0.02}{#1}}
\newcommand{\VariableTok}[1]{\textcolor[rgb]{0.00,0.00,0.00}{#1}}
\newcommand{\VerbatimStringTok}[1]{\textcolor[rgb]{0.31,0.60,0.02}{#1}}
\newcommand{\WarningTok}[1]{\textcolor[rgb]{0.56,0.35,0.01}{\textbf{\textit{#1}}}}
\setlength{\emergencystretch}{3em} % prevent overfull lines
\providecommand{\tightlist}{%
  \setlength{\itemsep}{0pt}\setlength{\parskip}{0pt}}
%\documentclass[unicode,10pt]{beamer}% 'unicode'が必要
\usepackage{luatexja}% 日本語を使う場合は必須
%\usepackage[japanese]{babel}	
%\usefonttheme[onlymath]{serif}
%\usepackage[T1]{fontenc} %これを使うと、ギリシャ文字が出ない
%\usepackage{textcomp}
%\usepackage[scale=1.0]{tgheros} % Sans serif font
%\usepackage[scaled]{beramono} % Monospace font
%\usepackage{luatexja-otf}
%\renewcommand{\kanjifamilydefault}{\gtdefault}
%\usepackage[haranoaji]{luatexja-preset}

% スライド番号の表示 %%%%%%%%%%%%%%%%%%%
\makeatletter
\setbeamertemplate{footline}{%
  \leavevmode%
  \hbox{%
  \begin{beamercolorbox}[wd=.5\paperwidth,ht=2.25ex,dp=1ex,right]{date in head/foot}%
    \insertframenumber{} \hspace*{2ex} % new version without total frames
  \end{beamercolorbox}}%
  \vskip0pt%
}
\makeatother
% スライド番号の表示 %%%%%%%%%%%%%%%%%%%
\usepackage{bookmark}
\IfFileExists{xurl.sty}{\usepackage{xurl}}{} % add URL line breaks if available
\urlstyle{same}
\hypersetup{
  hidelinks,
  pdfcreator={LaTeX via pandoc}}

\title{第7章: 不確実性の数値化}
\subtitle{Elena Llaudet \& Kosuke Imai.\\
Data Analysis for Social Science: A Friendly and Practical
Introduction.}
\author{}
\date{\vspace{-2.5em}2026-03-09}

\begin{document}
\frame{\titlepage}

\section{7.1
なぜ不確実性を測るのか？}\label{ux306aux305cux4e0dux78baux5b9fux6027ux3092ux6e2cux308bux306eux304b}

\begin{frame}{推測統計学の課題}
\phantomsection\label{ux63a8ux6e2cux7d71ux8a08ux5b66ux306eux8ab2ux984c}
\begin{itemize}
\tightlist
\item
  我々が手にするデータは通常、母集団からの\textbf{標本} (Sample)
  に過ぎない。
\item
  標本から計算された数値（標本平均など）は、抽出のたびに変動する。
\item
  \textbf{目標}:
  標本から得られた推定値が、母集団の真の値からどれくらい離れている可能性があるか（不確実性）を数値化する。
\item
  主要な道具：

  \begin{enumerate}
  \tightlist
  \item
    \textbf{信頼区間} (Confidence Intervals)
  \item
    \textbf{仮説検定} (Hypothesis Testing)
  \end{enumerate}
\end{itemize}
\end{frame}

\section{7.2 信頼区間}\label{ux4fe1ux983cux533aux9593}

\begin{frame}[fragile]{標本平均の95\%信頼区間}
\phantomsection\label{ux6a19ux672cux5e73ux5747ux306e95ux4fe1ux983cux533aux9593}
\begin{itemize}
\tightlist
\item
  標本平均を中心として、母平均が含まれる可能性が高い範囲。
\item
  数式（大標本の場合）: \(\bar{X} \pm 1.96 \times SE\)

  \begin{itemize}
  \tightlist
  \item
    \(SE\) (標準誤差)
    \(= \frac{\sigma}{\sqrt{n}} \approx \frac{S}{\sqrt{n}}\)
  \end{itemize}
\end{itemize}

\begin{Shaded}
\begin{Highlighting}[]
\NormalTok{bes }\OtherTok{\textless{}{-}} \FunctionTok{na.omit}\NormalTok{(}\FunctionTok{read.csv}\NormalTok{(}\StringTok{"BES.csv"}\NormalTok{))}
\NormalTok{n }\OtherTok{\textless{}{-}} \FunctionTok{nrow}\NormalTok{(bes)}
\NormalTok{sample\_mean }\OtherTok{\textless{}{-}} \FunctionTok{mean}\NormalTok{(bes}\SpecialCharTok{$}\NormalTok{leave)}
\NormalTok{se }\OtherTok{\textless{}{-}} \FunctionTok{sqrt}\NormalTok{(}\FunctionTok{var}\NormalTok{(bes}\SpecialCharTok{$}\NormalTok{leave)}\SpecialCharTok{/}\NormalTok{n)}

\CommentTok{\# 95\%信頼区間}
\FunctionTok{c}\NormalTok{(sample\_mean }\SpecialCharTok{{-}} \FloatTok{1.96} \SpecialCharTok{*}\NormalTok{ se, sample\_mean }\SpecialCharTok{+} \FloatTok{1.96} \SpecialCharTok{*}\NormalTok{ se)}
\end{Highlighting}
\end{Shaded}

\begin{verbatim}
## [1] 0.4657127 0.4780655
\end{verbatim}
\end{frame}

\begin{frame}[fragile]{平均の差の信頼区間}
\phantomsection\label{ux5e73ux5747ux306eux5deeux306eux4fe1ux983cux533aux9593}
\begin{itemize}
\tightlist
\item
  処置効果の推定値（平均の差）の不確実性を評価する。
\end{itemize}

\footnotesize

\begin{Shaded}
\begin{Highlighting}[]
\NormalTok{star }\OtherTok{\textless{}{-}} \FunctionTok{read.csv}\NormalTok{(}\StringTok{"STAR.csv"}\NormalTok{)}
\NormalTok{star}\SpecialCharTok{$}\NormalTok{small }\OtherTok{\textless{}{-}} \FunctionTok{ifelse}\NormalTok{(star}\SpecialCharTok{$}\NormalTok{classtype }\SpecialCharTok{==} \StringTok{"small"}\NormalTok{, }\DecValTok{1}\NormalTok{, }\DecValTok{0}\NormalTok{)}
\NormalTok{treatment }\OtherTok{\textless{}{-}}\NormalTok{ star}\SpecialCharTok{$}\NormalTok{reading[star}\SpecialCharTok{$}\NormalTok{small }\SpecialCharTok{==} \DecValTok{1}\NormalTok{]}
\NormalTok{control }\OtherTok{\textless{}{-}}\NormalTok{ star}\SpecialCharTok{$}\NormalTok{reading[star}\SpecialCharTok{$}\NormalTok{small }\SpecialCharTok{==} \DecValTok{0}\NormalTok{]}

\NormalTok{diff\_means }\OtherTok{\textless{}{-}} \FunctionTok{mean}\NormalTok{(treatment) }\SpecialCharTok{{-}} \FunctionTok{mean}\NormalTok{(control)}
\NormalTok{se\_diff }\OtherTok{\textless{}{-}} \FunctionTok{sqrt}\NormalTok{(}\FunctionTok{var}\NormalTok{(treatment)}\SpecialCharTok{/}\FunctionTok{length}\NormalTok{(treatment) }
                \SpecialCharTok{+} \FunctionTok{var}\NormalTok{(control)}\SpecialCharTok{/}\FunctionTok{length}\NormalTok{(control))}

\CommentTok{\# 95\%信頼区間}
\FunctionTok{c}\NormalTok{(diff\_means }\SpecialCharTok{{-}} \FloatTok{1.96} \SpecialCharTok{*}\NormalTok{ se\_diff, }
\NormalTok{  diff\_means }\SpecialCharTok{+} \FloatTok{1.96} \SpecialCharTok{*}\NormalTok{ se\_diff)}
\end{Highlighting}
\end{Shaded}

\begin{verbatim}
## [1]  3.167621 11.253473
\end{verbatim}

\normalsize
\end{frame}

\section{7.3 仮説検定}\label{ux4eeeux8aacux691cux5b9a}

\begin{frame}{帰無仮説と対立仮説}
\phantomsection\label{ux5e30ux7121ux4eeeux8aacux3068ux5bfeux7acbux4eeeux8aac}
\begin{itemize}
\tightlist
\item
  \textbf{帰無仮説 (\(H_0\))}:
  「効果はない（差はゼロである）」という仮定。
\item
  \textbf{対立仮説 (\(H_1\))}:
  「効果がある（差はゼロではない）」という主張。
\item
  \textbf{p値} (p-value):
  帰無仮説が正しいと仮定したとき、観測された値（またはそれ以上に極端な値）が得られる確率。

  \begin{itemize}
  \tightlist
  \item
    p値 \textless{} 0.05
    であれば、帰無仮説を棄却し、「統計的に有意」と判断する。
  \end{itemize}
\end{itemize}
\end{frame}

\begin{frame}[fragile]{平均の差の検定 (Z検定/t検定)}
\phantomsection\label{ux5e73ux5747ux306eux5deeux306eux691cux5b9a-zux691cux5b9atux691cux5b9a}
\begin{Shaded}
\begin{Highlighting}[]
\CommentTok{\# テスト統計量 (Z{-}score)}
\NormalTok{z\_obs }\OtherTok{\textless{}{-}}\NormalTok{ diff\_means }\SpecialCharTok{/}\NormalTok{ se\_diff}
\NormalTok{z\_obs}
\end{Highlighting}
\end{Shaded}

\begin{verbatim}
## [1] 3.495654
\end{verbatim}

\begin{Shaded}
\begin{Highlighting}[]
\CommentTok{\# p値の計算}
\DecValTok{2} \SpecialCharTok{*} \FunctionTok{pnorm}\NormalTok{(}\SpecialCharTok{{-}}\FunctionTok{abs}\NormalTok{(z\_obs))}
\end{Highlighting}
\end{Shaded}

\begin{verbatim}
## [1] 0.0004729011
\end{verbatim}

\begin{itemize}
\tightlist
\item
  p値が非常に小さいため、少人数学級の効果は統計的に有意である。
\end{itemize}
\end{frame}

\section{7.4
回帰係数の有意性}\label{ux56deux5e30ux4fc2ux6570ux306eux6709ux610fux6027}

\begin{frame}[fragile]{回帰分析の結果の解釈}
\phantomsection\label{ux56deux5e30ux5206ux6790ux306eux7d50ux679cux306eux89e3ux91c8}
\begin{itemize}
\tightlist
\item
  \texttt{summary()} で各係数の標準誤差、t統計量、p値を確認できる。
\end{itemize}

\footnotesize

\begin{Shaded}
\begin{Highlighting}[]
\NormalTok{uas }\OtherTok{\textless{}{-}} \FunctionTok{read.csv}\NormalTok{(}\StringTok{"UA\_survey.csv"}\NormalTok{)}
\NormalTok{fit }\OtherTok{\textless{}{-}} \FunctionTok{lm}\NormalTok{(pro\_russian\_vote }\SpecialCharTok{\textasciitilde{}}\NormalTok{ russian\_tv }\SpecialCharTok{+}\NormalTok{ within\_25km, }
          \AttributeTok{data =}\NormalTok{ uas)}
\FunctionTok{summary}\NormalTok{(fit)}\SpecialCharTok{$}\NormalTok{coefficients}
\end{Highlighting}
\end{Shaded}

\begin{verbatim}
##               Estimate Std. Error   t value     Pr(>|t|)
## (Intercept)  0.1959065 0.03457823  5.665602 3.032123e-08
## russian_tv   0.2875934 0.07652437  3.758194 2.000236e-04
## within_25km -0.2080644 0.07681051 -2.708802 7.079846e-03
\end{verbatim}

\normalsize

\begin{itemize}
\tightlist
\item
  \texttt{Pr(\textgreater{}\textbar{}t\textbar{})} 列が 0.05
  未満であれば、その変数の影響は統計的に有意である。
\end{itemize}
\end{frame}

\section{7.5 まとめ}\label{ux307eux3068ux3081}

\begin{frame}{第7章のまとめ}
\phantomsection\label{ux7b2c7ux7ae0ux306eux307eux3068ux3081}
\begin{itemize}
\tightlist
\item
  \textbf{不確実性}を無視してデータから結論を出すことはできない。
\item
  \textbf{信頼区間}は、推定値の「幅」を提示する。
\item
  \textbf{仮説検定}は、観測された結果が単なる偶然（サンプリングの偏り）で説明できるかどうかを判断する。
\item
  \textbf{p値 \textless{} 0.05}
  は慣習的な基準であるが、文脈に応じた解釈が重要である。
\end{itemize}
\end{frame}

\end{document}
